\documentclass[11pt]{article}
\usepackage{amsmath, amssymb, bm}
\usepackage{geometry}
\geometry{margin=1in}

\title{miniAV: Reactive Obstacle Navigation Using Proximity-Based State Estimation}
\author{Subhadeep Roy}
\date{}

\begin{document}

\maketitle

\section{Overview}

We design a reactive autonomous navigation system using three proximity sensors:
\[
S_k = 
\begin{bmatrix}
p_L \\
p_R \\
p_B
\end{bmatrix}_k
\]

The system consists of:

\begin{enumerate}
    \item A constant-velocity (CV) Kalman filter estimating proximity and approach rate.
    \item Latency-compensated prediction of future proximity.
    \item A state-space feedback controller producing linear and angular velocity commands.
\end{enumerate}

The goal is robust obstacle avoidance without explicit geometric reconstruction.

\section{1. Constant-Velocity State Estimation}

For each sensor $i \in \{L, R, B\}$, define state:

\[
x_k^{(i)} =
\begin{bmatrix}
p_k^{(i)} \\
\dot{p}_k^{(i)}
\end{bmatrix}
\]

State transition (constant velocity model):

\[
x_{k+1}^{(i)} =
\underbrace{
\begin{bmatrix}
1 & \Delta t \\
0 & 1
\end{bmatrix}
}_{A}
x_k^{(i)} + w_k^{(i)}
\]

Measurement model:

\[
z_k^{(i)} =
\underbrace{
\begin{bmatrix}
1 & 0
\end{bmatrix}
}_{H}
x_k^{(i)} + v_k^{(i)}
\]

where:

\begin{itemize}
    \item $w_k^{(i)} \sim \mathcal{N}(0, Q)$ is process noise
    \item $v_k^{(i)} \sim \mathcal{N}(0, R)$ is measurement noise
\end{itemize}

Kalman filter recursion is applied independently to each channel.

The estimator provides:

\[
\hat{p}_k^{(i)}, \quad \hat{\dot{p}}_k^{(i)}
\]

These represent smoothed proximity and estimated approach rate.

\section{2. Predict-Ahead Latency Compensation}

Let $\tau$ denote total control loop latency.

We predict future proximity using:

\[
\hat{p}_{k+\tau}^{(i)} =
\hat{p}_k^{(i)} + \tau \hat{\dot{p}}_k^{(i)}
\]

Define:

\[
p_L = \hat{p}_{k+\tau}^{(L)}, \quad
p_R = \hat{p}_{k+\tau}^{(R)}, \quad
p_B = \hat{p}_{k+\tau}^{(B)}
\]

Forward symmetric and differential components:

\[
p_{\text{sym}} = \frac{p_L + p_R}{2}
\]

\[
p_{\text{diff}} = p_R - p_L
\]

\[
p_{\max} = \max(p_L, p_R)
\]

These predicted values are used for control.

\section{3. State-Space Control Law}

Define control input:

\[
u_k =
\begin{bmatrix}
v_k \\
\omega_k
\end{bmatrix}
\]

where:
\begin{itemize}
    \item $v_k$ = linear velocity
    \item $\omega_k$ = angular velocity
\end{itemize}

\subsection*{Forward Velocity Law}

\[
v_k =
v_{\min} +
(v_{\max} - v_{\min})
\cdot
\text{clip}
\left(
1 - \frac{p_{\text{sym}}}{T_{\text{slow}}},
0, 1
\right)
\]

\subsection*{Steering Law}

\[
\omega_k = k_\omega \cdot p_{\text{diff}}
\]

\subsection*{Panic Reflex}

If:

\[
p_{\max} > T_{\text{panic}}
\]

then override:

\[
v_k = -v_{\text{escape}}
\]
\[
\omega_k = 
\begin{cases}
+\omega_{\text{escape}}, & p_L > p_R \\
-\omega_{\text{escape}}, & p_R > p_L
\end{cases}
\]

\subsection*{Rate Limiting}

To prevent oscillation:

\[
u_k = u_{k-1} + \text{clip}(u_k^\ast - u_{k-1}, -\Delta u_{\max}, \Delta u_{\max})
\]

where $u_k^\ast$ is raw command.

\section{4. Calibration Plan}

\subsection*{4.1 Sensor Noise Estimation}

Place robot stationary facing open space.

Collect samples:

\[
p_k^{(i)}
\]

Compute variance:

\[
R^{(i)} = \text{Var}(p_k^{(i)})
\]

\subsection*{4.2 Process Noise Tuning}

Drive robot at constant speed toward wall.

Adjust $Q$ such that:

\begin{itemize}
    \item $\hat{p}_k$ tracks smoothly
    \item $\hat{\dot{p}}_k$ responds quickly to approach
\end{itemize}

\subsection*{4.3 Threshold Selection}

\begin{itemize}
    \item $T_{\text{slow}}$: proximity at which slowdown begins.
    \item $T_{\text{panic}}$: proximity corresponding to collision risk.
\end{itemize}

These are empirically selected by measuring readings at known distances.

\subsection*{4.4 Latency Measurement}

Log timestamps of:

\begin{itemize}
    \item Command transmission
    \item Sensor notification reception
\end{itemize}

Estimate mean round-trip delay $\tau$ and use in prediction step.

\section{Conclusion}

The miniAV system combines:

\begin{itemize}
    \item Linear Kalman filtering of proximity signals,
    \item Short-horizon prediction for latency compensation,
    \item Nonlinear feedback control with safety overrides.
\end{itemize}

This avoids explicit geometric reconstruction while enabling stable
autonomous navigation in cluttered indoor environments.

\end{document}